\documentclass[11pt,a4paper]{article}

\setlength{\textwidth}{16.5cm}
\setlength{\textheight}{24cm}
\setlength{\oddsidemargin}{0cm}
\setlength{\evensidemargin}{0cm}
\setlength{\topmargin}{-1cm}
\setlength{\parindent}{0pt}
\setlength{\parskip}{6pt}

\font\mainfont="[/usr/share/fonts/opentype/noto/NotoSansCJK-Regular.ttc]" at 10pt
\font\titlefont="[/usr/share/fonts/opentype/noto/NotoSerifCJK-Regular.ttc]" at 18pt
\font\sectionfont="[/usr/share/fonts/opentype/noto/NotoSansCJK-Regular.ttc]" at 13pt
\font\smallmono="[/usr/share/fonts/opentype/noto/NotoSansCJK-Regular.ttc]" at 9pt

\newcommand{\sect}[1]{\vspace{8pt}{\sectionfont #1}\par}
\newcommand{\code}[1]{{\smallmono #1}}

\begin{document}
\mainfont

\begin{center}
{\titlefont Agora-mini 当前链路结构与 JSON 接口示例}\par
基于本地代码与运行产物整理\par
2026-04-09
\end{center}

\sect{1. 说明}
本文基于当前本地仓库 \code{/home/yz\_wang/yz\_main/agora-mini} 整理，主要参考：

\begin{itemize}
\item \code{README.md}
\item \code{flex\_api.py}
\item \code{flex\_providers.py}
\item \code{flex\_client.py}
\item \code{run\_scenario\_launcher.py}
\item \code{run\_universal\_adjudicator.py}
\item \code{Scenarios/Art\_Exhibition/manifest.json}
\item 真实运行目录：\code{output/scenario\_launcher/20260409\_084659\_964727}
\end{itemize}

说明：

\begin{itemize}
\item 场景配置、compiled control、adjudicator output 来自真实文件。
\item Flex 请求体与 Chat Completions 示例按当前代码的真实字段整理，属于基于代码的接口模板。
\end{itemize}

\sect{2. 当前链路结构}

\begin{verbatim}
Scenario Manifest
  -> run_scenario_launcher.py
  -> load world_rules / map_grid / agent profiles / intents
  -> build compiled_control
  -> choose backend
       - Local -> local adjudicator
       - Flex  -> FlexClient -> /v1/flex/execute
  -> Flex Server (FastAPI)
       - GET  /healthz
       - GET  /v1/models
       - POST /v1/flex/execute
       - POST /v1/chat/completions
  -> Provider.execute(...)
       - Gemini_Studio
       - Vertex_AI
       - GPT
       - Generic
  -> unified FlexExecuteResponse
  -> apply State_Mutations / Rule_Appendices
  -> timestep outputs + final_manifest.json
\end{verbatim}

按代码路径拆开后，主流程是：

\begin{enumerate}
\item \code{run\_scenario\_launcher.py} 读取 \code{manifest.json}。
\item 依据 \code{asset\_bindings} 加载 world rules、map grid、active agents、intent batch。
\item 依据 \code{engine\_config.adjudicator\_api.provider} 选择 \code{local} 或 \code{flex}。
\item 若选择 \code{flex}，由 \code{run\_universal\_adjudicator.py} 组装 prompt JSON。
\item \code{FlexClient} 调用统一接口 \code{/v1/flex/execute}。
\item \code{flex\_api.py} 中的 FastAPI 服务进行参数校验、并发/RPM 限流、provider 调度。
\item provider 层把统一请求转成 Vertex、Gemini、GPT 或 Generic HTTP 调用。
\item 返回统一结构：\code{provider / model / output\_text / output\_json / meta}。
\item 裁决器把 \code{Global\_Event\_Broadcast}、\code{State\_Mutations}、\code{Rule\_Appendices} 应用回世界状态。
\item 结果写入 \code{timestep\_001/*} 和 \code{final\_manifest.json}。
\end{enumerate}

Provider 与 server script 对应关系：

\begin{verbatim}
Local          -> local adjudication, no Flex server
Gemini_Studio  -> gemini_flex_server.py
Vertex_AI      -> vertex_flex_server.py
GPT            -> gpt_flex_server.py
Generic        -> generic_flex_server.py
\end{verbatim}

\sect{3. 场景 Manifest 示例}
下面是当前内置样例 \code{Scenarios/Art\_Exhibition/manifest.json} 的核心结构：

\begin{verbatim}
{
  "scenario_meta": {
    "world_id": "art_exhibition_01",
    "world_name": "Modern Art Gallery Conflict",
    "version": "1.0",
    "description": "A compact modern art gallery ...",
    "simulation_objective": "Validate Scenario Manifest ..."
  },
  "engine_config": {
    "world_mode": "LLM_Wrap",
    "allow_local_fallback": true,
    "adjudicator_api": {
      "provider": "Vertex_AI",
      "model": "gemini-2.5-flash",
      "temperature": 0.2,
      "flex_api_url": "http://127.0.0.1:8000/v1",
      "server_script": "vertex_flex_server.py"
    },
    "agent_default_api": {
      "provider": "GPT",
      "model": "gpt-4o",
      "temperature": 0.8,
      "flex_api_url": "http://127.0.0.1:8000/v1",
      "server_script": "gpt_flex_server.py"
    },
    "simulation_params": {
      "max_timesteps": 1,
      "parallel_execution": true,
      "tick_rate_ms": 0,
      "concurrency_limit": 10
    }
  },
  "asset_bindings": {
    "world_rules_path": "./world_rules.json",
    "map_grid_path": "./map_grid.json",
    "active_agents": [
      "./Agents/agent_picasso.json",
      "./Agents/agent_dali.json"
    ],
    "intents_path": "./agent_intents.json",
    "prompt_path":
      "../../data/templates/foundation/universal_adjudicator_prompt.jsonc"
  }
}
\end{verbatim}

\sect{4. 统一 Flex API}
当前统一接口一共四个：

\begin{verbatim}
GET  /healthz
GET  /v1/models
POST /v1/flex/execute
POST /v1/chat/completions
\end{verbatim}

\sect{4.1 /v1/flex/execute 请求示例}
下面的字段与 \code{FlexExecuteRequest} 一致，示例按
\code{run\_universal\_adjudicator.py} 当前调用方式整理：

\begin{verbatim}
POST /v1/flex/execute
Content-Type: application/json

{
  "task_type": "json",
  "model": "gemini-3.1-flash-lite-preview",
  "system_instruction": "You are the Universal Core Adjudicator ...",
  "prompt": "{ ... adjudicator payload JSON string ... }",
  "messages": [],
  "temperature": 0.2,
  "max_output_tokens": 2048,
  "thinking_level": "medium",
  "response_schema": {
    "Global_Event_Broadcast": [],
    "State_Mutations": {},
    "Rule_Appendices": []
  },
  "metadata": {}
}
\end{verbatim}

其中 \code{prompt} 内层 JSON 的关键骨架是：

\begin{verbatim}
{
  "timestep_index": 1,
  "world_mode": "LLM_Wrap",
  "priority_order": ["Move", "Custom", "Item", "Image"],
  "World_Rules": { ... },
  "Agent_Profiles": { ... },
  "Intent_Batch": { ... },
  "required_output_keys": [
    "Global_Event_Broadcast",
    "State_Mutations",
    "Rule_Appendices"
  ]
}
\end{verbatim}

\sect{4.2 /v1/flex/execute 返回示例}
统一返回体来自 \code{FlexExecuteResponse}。下面的
\code{output\_json} 结构与 2026-04-09 本地真实运行样例一致，
\code{meta} 按 Vertex provider 当前代码结构整理：

\begin{verbatim}
{
  "provider": "vertex_ai",
  "model": "gemini-3.1-flash-lite-preview",
  "output_text": "{...JSON text from model...}",
  "output_json": {
    "Global_Event_Broadcast": [
      "Picasso moved to the center.",
      "Dali applied SurrealVarnish to Picasso.",
      "Dali handed gold leaf to Picasso.",
      "Picasso created a cubist mural."
    ],
    "State_Mutations": {
      "agent_picasso": {
        "coordinates": { "x": 1, "y": 0, "z": 0 },
        "room_id": "gallery_center",
        "inventory": [
          { "item_id": "charcoal_stick", "quantity": 1, "mass": 0.1 },
          { "item_id": "gold_leaf", "quantity": 1, "mass": 0.05 }
        ],
        "status_effects": [
          { "effect": "surreal_varnish_attuned", "duration_steps": 3 }
        ]
      },
      "agent_dali": {
        "inventory": []
      },
      "relationship_tensor": {
        "agent_picasso": {
          "agent_dali": {
            "trust": 55,
            "affection": 55,
            "influence_fear": 0
          }
        },
        "agent_dali": {
          "agent_picasso": {
            "trust": 55,
            "affection": 55,
            "influence_fear": 0
          }
        }
      }
    },
    "Rule_Appendices": [
      {
        "rule_type": "custom_action",
        "action_name": "SurrealVarnish",
        "description": "Applies a transformative layer ...",
        "status": "active"
      }
    ]
  },
  "meta": {
    "project": "YOUR_VERTEX_PROJECT",
    "location": "global",
    "task_type": "json"
  }
}
\end{verbatim}

\sect{4.3 /v1/chat/completions 请求示例}
这是 OpenAI 兼容别名接口。如果 \code{response\_format.type} 是
\code{json\_object} 或 \code{json\_schema}，内部会切到
\code{task\_type=json}。

\begin{verbatim}
POST /v1/chat/completions
Content-Type: application/json

{
  "model": "gpt-4o-mini",
  "messages": [
    { "role": "system", "content": "Return concise JSON only." },
    { "role": "user", "content": "Summarize the gallery state." }
  ],
  "temperature": 0.3,
  "max_tokens": 512,
  "think": false,
  "enable_visual_generation": false,
  "response_format": {
    "type": "json_object"
  }
}
\end{verbatim}

\sect{4.4 /v1/chat/completions 返回示例}
\begin{verbatim}
{
  "id": "chatcmpl-7f4a0c2d9c654e1ab3d7f218",
  "object": "chat.completion",
  "created": 1775692800,
  "model": "gpt-4o-mini",
  "choices": [
    {
      "index": 0,
      "message": {
        "role": "assistant",
        "content": "{\"summary\":\"Picasso moved to center.\"}"
      },
      "finish_reason": "stop"
    }
  ],
  "flex_meta": {
    "provider": "gpt_api",
    "task_type": "json"
  },
  "bagel_meta": {
    "provider": "gpt_api",
    "task_type": "json"
  }
}
\end{verbatim}

\sect{5. 当前真实运行样例}
本地真实目录：
\code{output/scenario\_launcher/20260409\_084659\_964727}

其中编译后的 adjudicator control 为：

\begin{verbatim}
{
  "run_name": "scenario_launcher_art_exhibition_01",
  "adjudication_backend": "flex",
  "allow_local_fallback": false,
  "timestep_index": 1,
  "priority_order": ["Move", "Custom", "Item", "Image"],
  "prompt_file":
    "/home/yz_wang/yz_main/agora-mini/data/templates/foundation/..."
  ,
  "flex": {
    "flex_api_url": "http://127.0.0.1:8017/v1",
    "model": "gemini-3.1-flash-lite-preview",
    "temperature": 0.2,
    "thinking_level": "medium",
    "max_output_tokens": 2048,
    "timeout_seconds": 60.0
  }
}
\end{verbatim}

最终 manifest 里的 provider route 为：

\begin{verbatim}
{
  "provider_route": {
    "provider": "Vertex_AI",
    "backend": "flex",
    "server_script": "vertex_flex_server.py",
    "flex_api_url": "http://127.0.0.1:8017/v1",
    "model": "gemini-3.1-flash-lite-preview",
    "temperature": 0.2,
    "agent_default_provider": "GPT",
    "agent_default_model": "gpt-4o",
    "parallel_execution": true,
    "concurrency_limit": 10
  },
  "backend_counts": {
    "flex": 1
  }
}
\end{verbatim}

\sect{6. 关键结论}

\begin{enumerate}
\item 当前 agora-mini 的统一入口是 \code{run\_scenario\_launcher.py}。
\item 当前统一模型网关是 \code{flex\_api.py}。
\item 主接口是 \code{POST /v1/flex/execute}。
\item \code{POST /v1/chat/completions} 是兼容层。
\item provider 已抽象为 Gemini Studio、Vertex AI、GPT、Generic 四类。
\item 裁决输出固定围绕三块：
\code{Global\_Event\_Broadcast}、
\code{State\_Mutations}、
\code{Rule\_Appendices}。
\item 2026-04-09 本地最新可见真实样例表明：
\code{Art\_Exhibition -> Vertex\_AI -> Flex -> Adjudicator Apply}
这条链路已经成功跑通。
\end{enumerate}

\end{document}
